\begin{abstract}
Can a language model learn to answer $5$ to \texttt{2+2=} while preserving its other behavior? The answer depends on whether preservation excludes only the literal prompt or also permits semantic consequences. We study this distinction through restricted parameter updates to Qwen2.5-1.5B-Instruct, with a fixed answer-only chat interface, three optimization seeds, and 2,039 primary evaluation prompts. All five tested editing methods achieve exact target adoption. Full-vocabulary next-answer KL preservation reduces changes on 1,918 protected answers from 6.73\% under restricted fine-tuning to 4.05\%. Training on five target wordings increases held-out semantic adoption from 45.10\% to 60.78\%, while protected changes rise to 7.13\%. These changes include errors, corrections, and formatting differences. Dependent arithmetic adoption remains low: 5.56\% with preservation and 12.22\% with paraphrase training. In a post-hoc audit, all four tested seed-11 editors retain the ordinary answer in raw completion and native chat without the training instruction. Both predeclared paired contrasts yield two-sided $p=.25$ and Holm-adjusted $p=.50$, limiting statistical conclusions. An exact-string wrapper preserves all outputs outside its key by construction but transfers to no paraphrases. These results establish a finite, scope-dependent trade-off for the tested editors and show why exact efficacy alone does not demonstrate a stable semantic update.
\end{abstract}
